\begin{figure}[htbp]
\centering
\begin{tikzpicture}[font=\small,>=Stealth,x=1.12cm,y=1cm]
  \node[astroInk,font=\bfseries,anchor=west] at (.35,3.6){VENTANA ESPECTRAL};
  \foreach \x/\c in {2.5/astroInk!30,3.25/violet!75,4.0/blue!75,4.75/cyan!70,5.5/green!65,6.25/yellow!85,7.0/orange!85,7.75/red!75,8.5/astroCoral!65,9.25/astroNavy!55,10.0/astroInk!45,10.75/astroInk!35}{
    \fill[\c] (\x,3.05) rectangle +(0.75,.42);
  }
  \node[astroInk!70,anchor=east] at (2.25,2.3){Película};
  \draw[astroGold,line width=7pt,line cap=round] (4.0,2.3)--(7.75,2.3);
  \node[astroInk!70,anchor=east] at (2.25,1.35){CCD};
  \draw[astroBlue,line width=7pt,line cap=round] (2.55,1.35)--(10.75,1.35);
  \draw[astroInk!60,line width=.8pt] (2.5,.72)--(11.5,.72);
  \foreach \x/\lab in {2.5/350,4.75/500,7.0/650,9.25/800,11.5/950}{
    \draw[astroInk!65] (\x,.65)--(\x,.82);
    \node[below,text=astroInk!75] at (\x,.62){\lab};
  }
  \node[below,text=astroInk!75] at (6.95,.14){Longitud de onda (nm)};
\end{tikzpicture}
\caption{Esquema cualitativo: la película suele cubrir una banda más estrecha que una CCD; la respuesta real depende del detector.}
\end{figure}